\documentclass{article}
\input{../../lib.sty}
\allowdisplaybreaks

\title{\texttt{fn make\_noisy\_top\_k}}
\author{\MichaelShoemateAuthor}
\begin{document}
\maketitle
\contrib

Proves soundness of \rustdoc{measurements/fn}{make\_noisy\_top\_k}
in \asOfCommit{mod.rs}{f5bb719}.\\
\texttt{make\_noisy\_top\_k} returns a Measurement that
noisily selects the indices of the greatest scores from a vector of input scores.

\section{Hoare Triple}
\subsection*{Precondition}
\subsubsection*{Compiler-verified}
\begin{itemize}
    \item \texttt{MO} is a type with trait \rustdoc{measurements/make\_noisy\_top\_k/trait}{TopKMeasure}
    \item \texttt{TIA} (atomic input type) is a type with trait \rustdoc{traits/trait}{Number}
\end{itemize}

\subsubsection*{Caller-verified}
None

\subsection*{Pseudocode}
\label{sec:python-pseudocode}
\lstinputlisting[language=Python,firstline=2,escapechar=|]{./pseudocode/make_noisy_top_k.py}

\subsection*{Postcondition}

\begin{theorem}
    \label{postcondition}
    \validMeasurement{\texttt{input\_domain, input\_metric, output\_measure, k, scale, negate, MO, TIA}}{\texttt{make\_noisy\_top\_k}}
\end{theorem}

\begin{proof}[Proof of data-independent errors.]
    By the constructor checks and the postcondition of \rustdoc{measurements/noisy_top_k/fn}{noisy\_top\_k},
    the only source of error is due to entropy exhaustion, which could be data-dependent,
    due to differing number of expected random draws depending on the input dataset.

    The following privacy proof assumes an infallible source of independent uniform random bits.
    It does not establish a guarantee for observable entropy failures.
\end{proof}

\begin{proof}[Proof of privacy guarantee.]
    When \din\ is zero, by line~\ref{din-non-zero}, the privacy loss is zero, satisfying the postcondition.
    Otherwise when scale is zero, by line~\ref{scale-zero}, the privacy loss is infinite, also satisfying the postcondition.

    By the checks on lines~\ref{check-non-nan} and \ref{check-non-negative-scale},
    the preconditions for \rustdoc{measurements/noisy_top_k/fn}{noisy\_top\_k} are satisfied.
    Additionally by the checks on lines~\ref{din-non-neg}, \ref{din-non-zero} and \ref{scale-zero},
    the preconditions for \rustdoc{measurements/noisy_top_k/trait}{TopKMeasure}\texttt{privacy\_map} are satisfied.

    Clamping to finite bounds preserves the sign of each coordinate difference and does not increase its magnitude.
    Thus the range distance after clamping and optional negation is bounded by $\din$ for monotonic scores and $2\din$ otherwise, as computed with upward rounding in the privacy map.
    For a finite input distance, the score vectors have equal length.
    For any fixed transcript of selected indices, both vectors have the same coordinates removed,
    so the range distance of the remaining scores cannot increase.
    Thus the postcondition of \rustdoc{measurements/noisy_top_k/trait}{TopKMeasure} applies at every selection.
    Adaptive composition of the $\min(k, |x|)$ selections gives the bound
    on line~\ref{fn-privacy-map-call}.
\end{proof}

\end{document}
